% ============================================================
%  Detailed Description of the 2D Boussinesq AMOC Model
% ============================================================
\documentclass[12pt, a4paper]{article}

\usepackage[utf8]{inputenc}
\usepackage[T1]{fontenc}
\usepackage{amsmath, amssymb}
\usepackage{graphicx}
\usepackage{booktabs}
\usepackage{natbib}
\usepackage{hyperref}
\usepackage[margin=2.5cm]{geometry}
\usepackage{setspace}
\usepackage{siunitx}
\usepackage{cleveref}

\onehalfspacing

\renewcommand{\Pr}{\mathrm{Pr}}
\newcommand{\Le}{\mathrm{Le}}
\newcommand{\Ra}{\mathrm{Ra}}
\newcommand{\psu}{\,\text{psu}}
\newcommand{\CO}{CO\ensuremath{_2}}
\newcommand{\dint}{\mathrm{d}}
\newcommand{\pd}[2]{\frac{\partial #1}{\partial #2}}
\newcommand{\lap}{\nabla^2}

% ============================================================
\begin{document}

\title{Equations and Physical Interpretation of the
       Two-Dimensional Boussinesq AMOC Model}

\author{Andreas Morr}
\date{\today}
\maketitle

\begin{abstract}
This document provides a self-contained derivation and physical interpretation
of the two-dimensional latitude-depth Boussinesq model of the Atlantic
Meridional Overturning Circulation (AMOC) used in the AMOCResilience project.
The model is based on \citet{Soons2025Boussinesq} and encapsulates thermohaline
convection driven by prescribed surface temperature and salinity fluxes in an
overturning-relevant latitude-depth plane. We describe the governing equations
for temperature, salinity, vorticity, and streamfunction; the boundary
conditions; the surface forcing functions; the numerical discretisation; and
the Arctic amplification parameterisation used to sweep the model across a
range of \CO{}-equivalent forcing levels.
\end{abstract}

\tableofcontents
\newpage

% ============================================================
\section{Physical Setup and Non-Dimensionalisation}
\label{sec:setup}
% ============================================================

The model considers a two-dimensional, meridional cross-section of the ocean
at latitude $x \in [0, a]$ and depth $z \in [0, 1]$, where $a = 5$ is the
non-dimensional domain length corresponding to the full pole-to-pole extent
of the Atlantic basin ($x = 0$ is the South Pole, $x = a/2 = 2.5$ is the
equator, and $x = a = 5$ is the North Pole), and $z = 0$ is the ocean floor
and $z = 1$ is the surface. The model is formulated under the Boussinesq approximation, in which
density variations are neglected except in the buoyancy term. The flow is
assumed to be two-dimensional and incompressible.

Two scalar fields evolve dynamically: temperature $T(x,z,t)$ and salinity
$S(x,z,t)$. The velocity field is derived from a streamfunction $\psi(x,z,t)$
via
\begin{equation}
  u = \pd{\psi}{z}, \qquad w = -\pd{\psi}{x},
  \label{eq:vel}
\end{equation}
where $u$ is the horizontal (meridional) velocity component and $w$ is the
vertical velocity. The flow is characterised by the vorticity
$\omega = \partial_x w - \partial_z u = -\lap\psi$.

The non-dimensional parameters governing the system are:
\begin{itemize}
  \item the \textbf{Prandtl number} $\Pr = \nu/\kappa_T$, the ratio of
        kinematic viscosity to thermal diffusivity;
  \item the \textbf{Lewis number} $\Le = \kappa_T/\kappa_S$, the ratio of
        thermal to haline diffusivity;
  \item the \textbf{Rayleigh number} $\Ra$, which measures the strength of
        buoyancy-driven convection relative to diffusion;
  \item the \textbf{aspect ratio} $A = a$ (domain length in non-dimensional
        units).
\end{itemize}
The default parameter values used in all simulations reported in the
AMOCResilience paper are listed in \Cref{tab:params}.

\begin{table}[h]
  \centering
  \caption{Default non-dimensional model parameters.}
  \label{tab:params}
  \begin{tabular}{lll}
    \toprule
    Symbol & Value & Physical meaning \\
    \midrule
    $\Pr$        & $1$              & Prandtl number \\
    $\Le$        & $1$              & Lewis number (thermal/haline diffusivity ratio) \\
    $\Ra$        & $4\times10^4$    & Rayleigh number \\
    $a$          & $5$              & Non-dimensional domain length \\
    $\delta$     & $0.05$           & Surface boundary-layer thickness \\
    $\tau_T$     & $0.1$            & Thermal relaxation time scale \\
    $\tau_S$     & $1.0$            & Salinity forcing time scale \\
    $\beta$      & $0.1$            & Salinity-forcing asymmetry parameter \\
    $N_k$        & $7$              & Number of stochastic salinity modes \\
    $M,\,N$      & $20,\,40$        & Latitude, depth grid points \\
    $\Delta t$   & $10^{-3}$        & Time step \\
    \bottomrule
  \end{tabular}
\end{table}

% ============================================================
\section{Governing Equations}
\label{sec:equations}
% ============================================================

\subsection{Temperature}
\label{sec:temp}

Temperature evolves according to an advection-diffusion equation with surface
restoring:
\begin{equation}
  \pd{T}{t} + u\pd{T}{x} + w\pd{T}{z}
  = \lap T + \frac{h(z)}{\tau_T}\bigl(T_s(x) - T\bigr).
  \label{eq:T}
\end{equation}
The left-hand side is the material derivative of temperature. The first term
on the right-hand side, $\lap T = \partial_{xx}T + \partial_{zz}T$, represents
isotropic diffusion. The second term is a relaxation towards the prescribed
surface temperature profile $T_s(x)$, active only in a thin surface layer
characterised by the exponential decay function
\begin{equation}
  h(z) = \exp\!\left(-\frac{1-z}{\delta}\right),
  \label{eq:h}
\end{equation}
with $\delta = 0.05$. Since $h(z) \approx 1$ near the surface ($z \to 1$) and
decays rapidly to zero at depth, the forcing is effectively confined to the
uppermost portion of the water column, mimicking the role of the atmospheric
heat exchange. The time scale $\tau_T = 0.1$ controls the strength of this
restoring; a smaller $\tau_T$ enforces temperature more strongly towards $T_s$.

The surface temperature profile is prescribed as
\begin{equation}
  T_s(x) = \frac{1}{2}\!\left[\cos\!\left(\frac{2\pi x}{a} - \pi\right) + 1\right],
  \label{eq:Ts}
\end{equation}
This profile is symmetric around the equator: $T_s = 0$ at both poles
($x = 0$ and $x = a$, representing the cold Southern and Northern Poles,
respectively) and $T_s = 1$ at the equator ($x = a/2 = 2.5$), representing
the warm tropics. The equator-to-pole temperature contrast drives the
horizontal density gradient that sustains the overturning circulation, with
both hemispheres contributing symmetrically to the thermal forcing.

\subsection{Salinity}
\label{sec:salt}

Salinity obeys a similar advection-diffusion equation, but with a prescribed
surface flux rather than a relaxation, and with an additional stochastic
forcing term:
\begin{equation}
  \pd{S}{t} + u\pd{S}{x} + w\pd{S}{z}
  = \frac{1}{\Le}\lap S
    + \frac{h(z)}{\tau_S}\,S_s(x;\beta)
    + h(z)\,\sigma_S\,\xi(x,t).
  \label{eq:S}
\end{equation}
The diffusivity is rescaled by $1/\Le$ relative to the thermal equation. The
prescribed surface salinity flux $S_s(x;\beta)$ is
\begin{equation}
  S_s(x;\beta) =
    3.5\cos\!\left(\frac{2\pi x}{a} - \pi\right)
    - \beta\sin\!\left(\frac{\pi x}{a} - \frac{\pi}{2}\right),
  \label{eq:Ss}
\end{equation}
with asymmetry parameter $\beta = 0.1$. The dominant cosine term represents
the atmospheric freshwater cycle: the tropics receive a net salinity flux
(evaporation excess) while high latitudes receive a net freshwater flux
(precipitation excess). The sine term with amplitude $\beta$ introduces a
weak hemispheric asymmetry that breaks the symmetry of the cosine, slightly
enhancing the freshwater input in the northern high latitudes relative to the
southern high latitudes. This asymmetry is physically motivated by the greater
land area and different hydrological cycle of the Northern Hemisphere.

In contrast to temperature, which is restored to $T_s$, salinity is forced
by a fixed flux $S_s$ and does not relax to a reference profile. This is
physically appropriate: ocean salinity is set by freshwater fluxes from
precipitation, evaporation, and river runoff, not by direct exchange with an
atmospheric salinity reservoir. A global salinity correction is applied at
each time step to conserve domain-mean salinity exactly.

The stochastic term $h(z)\,\sigma_S\,\xi(x,t)$ represents unresolved
mesoscale eddy variability and atmospheric freshwater flux noise. It is
expanded in $N_k$ Fourier modes in latitude:
\begin{equation}
  h(z)\,\sigma_S\,\xi(x,t)
  = h(z)\sqrt{\frac{\varepsilon}{N_k}}
    \sum_{k=1}^{N_k}\bigl[
      \eta_k^c(t)\cos\!\tfrac{2\pi kx}{a}
      + \eta_k^s(t)\sin\!\tfrac{2\pi kx}{a}
    \bigr],
\end{equation}
where $\eta_k^c, \eta_k^s$ are independent standard Gaussian white-noise
increments and $\varepsilon$ is the noise variance. In the deterministic
resilience sweeps described in the AMOCResilience paper, the noise amplitude
is set to $\varepsilon = 0.05$ during the spinup phase and to zero during the
perturbation experiments to ensure deterministic basin mapping.

\subsection{Vorticity}
\label{sec:vort}

The vorticity equation is obtained by taking the curl of the momentum equation
under the Boussinesq approximation. Viscous dissipation and advection of
vorticity are balanced by the baroclinic production of vorticity due to the
horizontal density gradient:
\begin{equation}
  \pd{\omega}{t} + u\pd{\omega}{x} + w\pd{\omega}{z}
  = \Pr\,\lap\omega + \Pr\,\Ra\,\pd{}{x}(T - S).
  \label{eq:omega}
\end{equation}
The first term on the right-hand side is viscous diffusion of vorticity,
scaled by $\Pr$. The second term is the buoyancy torque: in the Boussinesq
approximation, the linearised equation of state gives density $\rho \propto
-(T - S)$ (assuming equal coefficients of thermal expansion and haline
contraction in non-dimensional units), so horizontal density gradients drive
vorticity production. A warmer or fresher column at the same latitude is
lighter, and the resulting meridional density contrast drives overturning.

The product $\Pr\cdot\Ra = 4\times10^4$ controls the ratio of advective to
diffusive vorticity transport and is the primary control on whether a
circulation develops at all. For $\Ra$ above a critical value, the baroclinic
forcing overcomes diffusion and sustains a mean overturning cell.

\subsection{Streamfunction and Incompressibility}
\label{sec:psi}

The streamfunction $\psi$ is related to vorticity through the two-dimensional
Poisson equation that follows from incompressibility and the definition
$\omega = -\lap\psi$:
\begin{equation}
  \lap\psi = -\omega.
  \label{eq:psi}
\end{equation}
This equation is solved at each time step given the updated vorticity field.
The streamfunction fully characterises the two-dimensional velocity field via
\Cref{eq:vel}. The AMOC strength in physical units is proportional to the
maximum value of $\psi$ in the interior; the maximum streamfunction $\psi_\mathrm{ref}$
at the unperturbed equilibrium provides the calibration anchor for converting
to Sverdrups (see \Cref{sec:scaling}).

% ============================================================
\section{Boundary Conditions}
\label{sec:bcs}
% ============================================================

No-flux and no-slip conditions are imposed on all lateral boundaries. For
temperature and salinity, the Neumann conditions $\partial_n T = \partial_n S = 0$
hold on all boundaries, meaning there is no diffusive flux through the walls.
The surface forcing enters through the interior relaxation and flux terms
(multiplied by $h(z)$) rather than through inhomogeneous boundary conditions,
which simplifies the numerical treatment.

For vorticity, free-slip conditions are applied on the lateral walls ($x = 0$
and $x = a$) and no-slip conditions on the top ($z = 1$) and bottom ($z = 0$)
boundaries. For the streamfunction, the Dirichlet condition $\psi = 0$ is
imposed on all boundaries, consistent with a closed domain with no net
overturning flux through any wall.

% ============================================================
\section{Depth-Stretched Grid}
\label{sec:grid}
% ============================================================

The depth coordinate is discretised on a non-uniform grid with $N + 1$ points
$z_j$ defined by the mapping
\begin{equation}
  z_j = \frac{1}{2} + \frac{\tanh\!\bigl(q(j/N - \tfrac{1}{2})\bigr)}{2\tanh(q/2)},
  \qquad j = 0, 1, \ldots, N,
  \label{eq:zgrid}
\end{equation}
with stretching parameter $q = 3$. This mapping concentrates grid points near
both the surface ($z \approx 1$) and the bottom ($z \approx 0$), where the
boundary layers and strong vertical gradients are located. The latitude
coordinate uses a uniform grid with $M + 1$ points $x_i = ia/(M)$ for
$i = 0,\ldots,M$.

All spatial derivatives are approximated by second-order finite differences.
The second-derivative operators $D_{xx}$ and $D_{zz}$ use Neumann boundary
conditions (the flux variants $F_{xx}$, $F_{zz}$ use Dirichlet conditions for
$\psi$ and $\omega$). First-derivative operators $D_x$ and $D_z$ are central
differences with zero-gradient corrections on the walls.

% ============================================================
\section{Time Stepping}
\label{sec:timestepping}
% ============================================================

Equations~\eqref{eq:T}, \eqref{eq:S}, and \eqref{eq:omega} are time-stepped
using an implicit-explicit (IMEX) scheme with time step $\Delta t = 10^{-3}$.
The diffusion and surface-relaxation terms are treated implicitly to avoid the
stability restriction of the diffusion CFL condition; the advection terms are
treated explicitly.

The implicit step for each variable leads to a generalised Sylvester equation
of the form
\begin{equation}
  A\,X + X\,B = C,
\end{equation}
where $X \in \mathbb{R}^{(M+1)\times(N+1)}$ is the updated field, $A$ and $B$
contain the implicit operators in $x$ and $z$ respectively, and $C$ contains
the explicit contributions and the field at the previous time step. This
separable structure arises because the diffusion operator is a sum of a
pure-$x$ and a pure-$z$ operator. For temperature, for example:
\begin{align}
  A_T &= I - \Delta t\,D_{xx}, \\
  B_T &= \frac{\Delta t}{\tau_T}H_z - \Delta t\,D_{zz},\\
  C_T &= T^n + \Delta t\,\bigl(-\psi^n_z\,\partial_x T^n
         + \psi^n_x\,\partial_z T^n + F_T\bigr),
\end{align}
where $H_z = \mathrm{diag}(h(z_j))$ is the diagonal surface-weight matrix
and $F_T$ is the surface temperature forcing. The Sylvester equations are
solved at each step using the Bartels--Stewart algorithm as implemented in
\texttt{scipy.linalg.solve\_sylvester}.

% ============================================================
\section{Arctic Amplification Parameterisation}
\label{sec:gamma}
% ============================================================

To sweep the model across a range of climate states corresponding to different
\CO{} concentrations, the surface temperature forcing is modified to mimic
Arctic amplification. The effective surface temperature profile is blended
between the standard pole-to-equator profile $T_s(x)$ and its spatial mean:
\begin{equation}
  T_s^\gamma(x) = (1-\gamma)\,T_s(x) + \gamma\,\overline{T_s},
  \label{eq:Tgamma}
\end{equation}
where $\overline{T_s} = \frac{1}{a}\int_0^a T_s(x)\,\mathrm{d}x = \tfrac{1}{2}$
is the domain-mean surface temperature. The parameter $\gamma \in [0,1]$
measures the degree of Arctic amplification: $\gamma = 0$ leaves the forcing
unchanged (pre-industrial-like), while $\gamma = 1$ eliminates the
equator-to-pole temperature gradient entirely, abolishing the thermal driving
of the overturning. Intermediate values reduce the meridional temperature
contrast, weakening the baroclinic torque in \Cref{eq:omega} and thereby
modifying the stability and strength of the AMOC.

A physically grounded mapping between $\gamma$ and atmospheric \CO{}
concentration $C$ is derived from the logarithmic forcing formula and an
estimate of Arctic ocean-surface amplification (see \texttt{co2\_gamma\_mapping.py}):
\begin{equation}
  \gamma(C) = \frac{(\alpha_\mathrm{sst}-1)\,\mathrm{ECS}}{\Delta T_0}
              \log_2\!\left(\frac{C}{C_0}\right)
  = 0.1\log_2\!\left(\frac{C}{280\,\text{ppm}}\right),
  \label{eq:co2gamma}
\end{equation}
using ECS $= 3.0\,^{\circ}\text{C}$ per \CO{} doubling, polar ocean-surface
amplification $\alpha_\mathrm{sst} = 2.0$, pre-industrial gradient
$\Delta T_0 = 30\,^{\circ}\text{C}$, and $C_0 = 280\,\text{ppm}$. The sweep
covers 15 uniformly spaced values of $\gamma$ in $[0, 0.14]$, corresponding
to \CO{} concentrations from $280\,\text{ppm}$ (pre-industrial) to approximately
$740\,\text{ppm}$ (well beyond a doubling of \CO{}).

The physical interpretation of \Cref{eq:Tgamma} is as follows. Under Arctic
amplification, warming in high latitudes outpaces warming in low latitudes,
reducing the pole-to-equator SST difference. This weakens the meridional
density gradient between the North Atlantic (cold, dense) and the tropics
(warm, light) that drives NADW formation. For the Boussinesq model, this
manifests as a reduction of the buoyancy torque $\Pr\cdot\Ra\,\partial_x(T-S)$
at the surface, which feeds into the vorticity budget and modifies the
equilibrium streamfunction strength. As $\gamma$ increases, the AMOC-on
equilibrium strengthens slightly (the reduced thermal gradient allows the
haline component to dominate) while the convergence time and characteristic
return time decrease, consistent with a more stable on-state at higher thermal
mixing.

% ============================================================
\section{Geographic Box Perturbations}
\label{sec:boxes}
% ============================================================

The resilience experiments apply coherent salinity perturbations to two
geographic regions of the model domain, analogous to the Wood-box geometry
used in the CLIMBER-X experiments. Each box is defined by a smooth mask
$m_\mathcal{B}(x, z)$ that localises the perturbation spatially while avoiding
numerical artefacts from sharp boundaries. The mask is the product of a
surface-depth profile and a tanh lateral transition:
\begin{equation}
  m_\mathcal{B}(x,z) = h(z)\cdot\tfrac{1}{2}\left[
    \tanh\!\left(\frac{x - x_\mathrm{lo}}{\ell}\right)
    - \tanh\!\left(\frac{x - x_\mathrm{hi}}{\ell}\right)
  \right]_+,
\end{equation}
normalised so that the mask-weighted box-mean equals one. The transition
half-width is $\ell = 0.25$ in non-dimensional latitude coordinates. The
surface-depth weighting $h(z)$ concentrates the perturbation in the upper
ocean where freshwater anomalies are observed.

Two boxes are used:
\begin{itemize}
  \item \textbf{North Atlantic box}: $x \in [130/36, 5]$, approximately $40{^\circ}$--$90{^\circ}\text{N}$;
  \item \textbf{Tropical Atlantic box}: $x \in [40/36, 130/36]$, approximately $50{^\circ}\text{S}$--$40{^\circ}\text{N}$.
\end{itemize}
A perturbation $(\Delta S_N, \Delta S_T)$ is applied by finding the field
amplitudes that achieve the target box-mean shifts simultaneously, including
cross-box effects. A uniform compensating offset is subtracted from the full
domain to conserve domain-mean salinity exactly. Perturbations are swept on a
$15\times15$ grid with both coordinates spanning $[-2, 2]\psu{}$.

% ============================================================
\section{Convergence and Physical Scaling}
\label{sec:scaling}
% ============================================================

After a salinity perturbation is applied, the model is integrated forward
until convergence. Convergence is declared when the Euclidean distance between
the current North/Tropical box-salinity vector $(S_N, S_T)$ and the
unperturbed reference $(S_N^*, S_T^*)$ falls below
$\varepsilon = 0.05\psu{} \approx 0.005$ model units. The mean convergence
time across all perturbations that return to the on-state provides the
principal resilience metric from this model.

Since all governing equations are non-dimensional, the output quantities are
converted to physically interpretable units for comparison with other models
in the synthesis figure:

\paragraph{AMOC strength.}
The overturning streamfunction $\psi$ is calibrated by anchoring the
pre-industrial ($\gamma = 0$) equilibrium value
$\psi_0 \approx 1.52\,\text{model units}$ to the observed RAPID array estimate
and palaeoclimate reconstructions of the pre-industrial AMOC strength of
$18\,\text{Sv}$, giving the linear conversion
\begin{equation}
  \psi_\mathrm{Sv} = \frac{18.0}{\psi_0}\,\psi
  \approx 11.9\,\text{Sv}\times\psi.
\end{equation}

\paragraph{Time scales.}
Convergence times and characteristic return times are reported in
non-dimensional model time units. A conversion factor of $100$ is applied,
mapping the pre-industrial mean convergence time ($\approx 1.49$ model units)
and characteristic return time ($\approx 0.47$ model units) to year-proxy
values of $\approx 149\,\text{yr}$ and $\approx 47\,\text{yr}$, respectively.
These values are broadly consistent with the centennial-scale relaxation times
found in full ocean GCMs. The scaling factor is fixed by a single pre-industrial
calibration point; all trends with $\gamma$ are unaffected by the choice of
scale.

\paragraph{Characteristic return time.}
The local relaxation time scale $\tau = -1/\mathrm{Re}(\lambda_\mathrm{max})$
is estimated from the dominant eigenvalue of the $2\times2$ linearised map
in the $(S_N, S_T)$ box-salinity subspace. Eight small isotropic perturbations
($\pm 0.5\psu{}$ along directions at $45{^\circ}$ increments) are applied to the
equilibrium state; the $2\times2$ matrix $M$ mapping initial to final
deviations after $N_\mathrm{steps}$ steps is estimated by least squares, and
the continuous-time generator $A = \log(M)/\Delta t$ is eigendecomposed to
extract $\tau$.

% ============================================================
\section{Summary of the Equation System}
\label{sec:summary}
% ============================================================

For reference, the full system of governing equations is collected below:
\begin{align}
  \pd{T}{t} + u\pd{T}{x} + w\pd{T}{z}
    &= \lap T + \frac{h(z)}{\tau_T}\bigl(T_s^\gamma(x) - T\bigr),
    \label{eq:full_T}\\[4pt]
  \pd{S}{t} + u\pd{S}{x} + w\pd{S}{z}
    &= \frac{1}{\Le}\lap S + \frac{h(z)}{\tau_S}S_s(x;\beta)
       + h(z)\,\sigma_S\,\xi(x,t),
    \label{eq:full_S}\\[4pt]
  \pd{\omega}{t} + u\pd{\omega}{x} + w\pd{\omega}{z}
    &= \Pr\,\lap\omega + \Pr\cdot\Ra\,\pd{}{x}(T - S),
    \label{eq:full_omega}\\[4pt]
  \lap\psi &= -\omega,
    \label{eq:full_psi}\\[4pt]
  u &= \pd{\psi}{z}, \qquad w = -\pd{\psi}{x},
    \label{eq:full_vel}
\end{align}
with surface temperature forcing
\begin{equation}
  T_s^\gamma(x) = (1-\gamma)\cdot\frac{1}{2}\!\left[
    \cos\!\!\left(\frac{2\pi x}{a} - \pi\right) + 1
  \right] + \frac{\gamma}{2},
\end{equation}
surface salinity forcing
\begin{equation}
  S_s(x;\beta) = 3.5\cos\!\!\left(\frac{2\pi x}{a} - \pi\right)
               - \beta\sin\!\!\left(\frac{\pi x}{a} - \frac{\pi}{2}\right),
\end{equation}
and surface boundary-layer profile $h(z) = \exp(-(1-z)/\delta)$.

% ============================================================
\bibliographystyle{apalike}
\begin{thebibliography}{1}

\bibitem[Soons et~al.(2025)]{Soons2025Boussinesq}
Soons, J., Grafke, T., and Dijkstra, H.~A. (2025).
\newblock Constructing efficient score functions for the rare event simulation
  of {AMOC} transitions.
\newblock \emph{Journal of Fluid Mechanics}, 10.1017/jfm.2025.248.

\end{thebibliography}

\end{document}
